\documentclass{article}
\usepackage{amsmath, amssymb, amsthm}

\setlength{\parindent}{0pt}

\newtheorem*{claim}{Claim}

\begin{document}

\begin{claim}
    There exist arbitrarily long runs of consecutive composite numbers.
\end{claim}

\begin{proof}
    Let $n \ge 2$ be arbitrary, and consider the $n-1$ consecutive integers
    \[
    n!+2,\; n!+3,\; \dots,\; n!+n
    \]
    For each integer $k$ with $2 \le k \le n$, $k$ is one of the factors of
    \[
    n! = \prod_{i=1}^{n} i
    \]
    so $k \mid n!$. Thus
    \[
    n! \equiv 0 \pmod{k}
    \]
    and therefore
    \[
    n! + k \equiv 0 + k \equiv 0 \pmod{k}
    \]
    Hence $k \mid (n!+k)$.

    \medskip

    Moreover,
    \[
    n!+k > k,
    \]
    since $n! \ge 2$ and $k \ge 2$. Therefore $k$ is a proper divisor of $n!+k$, so each integer 
    is composite.

    \medskip

    Since $n$ was arbitrary, there exist arbitrarily long runs of consecutive composite numbers.
\end{proof}

\end{document}
